\chapter{Introdu{\c c}\~ao}

%Neste capítulo são apresentadas as motivações para a realização da análise temporal de membranas comerciais no tratamento de água de produção via destilação por membranas, os objetivos do estudo e a estrutura da dissertação.

A extração de petróleo e gás natural gera como principal subproduto a água de produção, com geração estimada em aproximadamente 250 milhões de barris por dia (39,75 milhões de metros cúbicos)~\cite{w15162980, MembranePW_JSDEWES_2024}. Estima-se que, somente em 2018, a geração anual mundial de água de produção tenha atingido 6,8 bilhões de metros cúbicos, volume que tende a aumentar com o envelhecimento dos campos petrolíferos, podendo representar até 98~\% dos fluidos produzidos em poços maduros~\cite{w15234088,ALGHOUTI2019222}. 

Historicamente tratada como um rejeito da indústria petrolífera, a água de produção passou a ser reconhecida como um recurso de valor estratégico, especialmente devido ao seu potencial de reutilização em operações de reinjeção, recuperação avançada de petróleo e o aumento do rigor de regulamentações ambientais.

A água de produção é composta pela mistura entre água de formação, naturalmente presente nos reservatórios, e água de injeção, utilizada para manutenção da pressão do reservatório e mobilização do óleo residual. Sua composição varia de acordo com a geologia do reservatório e localização geográfica, podendo conter óleo disperso, sais dissolvidos, metais pesados, sólidos suspensos, compostos orgânicos e produtos químicos empregados nas operações de produção ~\cite{BEYER2020105155, ALGHOUTI2019222, w13020183}.

A crescente demanda por práticas ambientais sustentáveis e pela maximização da recuperação de hidrocarbonetos tem impulsionado o desenvolvimento de tecnologias capazes de viabilizar o reaproveitamento da água produzida. Nesse contexto, a reinjeção destaca-se como alternativa de elevada relevância técnica e econômica, uma vez que permite reduzir o consumo de água fresca, minimizar o descarte ambiental e aumentar o fator de recuperação de petróleo por meio da manutenção da pressão do reservatório e da mobilização de óleo remanescente.

Nesse contexto, o reaproveitamento da água de produção deixa de representar apenas uma alternativa para adequação ambiental, passando a desempenhar papel estratégico na indústria do petróleo. 

A reinjeção da água tratada contribui simultaneamente para atender regulamentações~\cite{OSPAR2001,EPA435,CONAMA3932007}, diminuir o consumo de água fresca e aumento do fator de recuperação de hidrocarbonetos, por meio da manutenção da pressão do reservatório e da mobilização de óleo residual. 

Entretanto, a viabilidade dessa prática está diretamente associada à qualidade da água empregada, uma vez que elevadas concentrações de óleo, sólidos suspensos e compostos incrustantes podem provocar danos ao reservatório e aos sistemas de injeção~\cite{BADER2007159, AKSTINAT2019124011}.


Diante dessa necessidade, diferentes etapas de tratamento têm sido empregadas na indústria do petróleo, abrangendo processos primários, secundários, terciários e de polimento~\cite{Amakiri2022ProducedWaterReview,Adepitan2026HybridPW,Rajbongshi2024ProducedWaterReview,Ahmadun2009PWTechnologies}. Entre as tecnologias emergentes para tratamento e reaproveitamento da água de produção, a destilação por membranas destaca-se pela elevada rejeição de sais, capacidade de operar com soluções de alta salinidade e potencial integração com fontes térmicas de baixa temperatura.

%A extração de petróleo e gás apresenta como principal subproduto a água de produção, com uma geração aproximada de 250 milhões de barris (39,75 milhões de metros cúbicos) por dia~\cite{w15162980, MembranePW_JSDEWES_2024}. Estima-se que no ano de 2018 houve uma geração anual mundial de 6,8 bilhões de metros cúbicos de água de produção, o que promove desafios econômicos e ambientais para as empresas do ramo petrolífero, uma vez que o volume de água de produção gerado aumenta com o envelhecimento do poço e pode representar até 98\% do material extraído~\cite{w15234088,ALGHOUTI2019222}. 

%Água de produção é o termo utilizado para o conjunto de águas oriundas do processo de extração, sendo composta por água de formação (água presente naturalmente nos poços de petróleo) e água de injeção (água utilizada para manutenção da pressão do poço e auxiliar na remoção do petróleo), apresentando características e contaminantes que variam em função da geologia do ponto de extração, da localização geográfica e da profundidade do poço~\cite{BEYER2020105155, ALGHOUTI2019222, w13020183}.

%Como principais constituintes da água de produção, é possível destacar óleo disperso, gases de hidrocarbonetos dissolvidos, sais inorgânicos, ácidos orgânicos, hidrocarbonetos aromáticos, cetonas, fenóis, metais pesados, materiais radioativos de ocorrência natural, partículas sólidas em suspensão, anti-incrustantes, anticorrosivos e biocidas~\cite{Ozgun, ALLEY201174}.  

%Uma vez extraída, a água de produção precisa ser tratada de acordo com sua destinação, seja para descarte ou reinjeção. Em caso de descarte, é necessário observar as regulamentações ambientais do local de extração, como a recomendação 2001/1 da convenção Oslo Paris, que determina uma concentração de 30~mg/L de óleo em água~\cite{OSPAR2001}; a agência de proteção ambiental dos Estados Unidos (USEPA) e a resolução número 393/2007 do Conselho Nacional do Meio Ambiente (CONAMA) determinam uma concentração média mensal inferior a 29~mg/L para óleos e graxas e um valor diário máximo de 42~mg/L~\cite{EPA435,CONAMA3932007}.

%Apesar das regulações impostas pelas organizações citadas anteriormente, o descarte contínuo de água de produção é um ponto de atenção, uma vez que é a principal fonte de hidrocarbonetos aromáticos policíclicos (PAHs) e petróleo em ambiente marinho, podendo ocasionar problemas agudos ou crônicos a biota marinha \cite{BEYER2020105155}. 

%No que diz respeito à água para reinjeção, a qualidade da água determina parâmetros de injeção como pressão e obstrução de poros. Nesse sentido, a análise química da água utilizada determina a dosagem adequada de anti-incrustantes e anti-corrosivos, além de atender aos requisitos recomendados de conter no máximo 10~mg/L de sólidos suspensos totais (TSS) e 42~mg/L de óleo em água. Outro ponto de atenção é a compatibilidade química entre a água a ser utilizada na reinjeção e a água de formação presente no poço, de modo a evitar reações que propiciem a precipitação de compostos inorgânicos e a formação de encrustações~\cite{BADER2007159, AKSTINAT2019124011}. 

%Mediante a necessidade de tratamento da água de produção, diferentes etapas de tratamento são empregadas para atender os requisitos das normas citadas anteriormente. De modo geral,  tais etapas de tratamento são descritas na literatura de acordo com a tecnologia empregada e o nível de contaminantes presentes, compondo quatro grupos principais: primário, secundário, terciário e polimento, de modo que a destilação por membranas pode estar presente no grupo terciário ou na etapa de polimento~\cite{Amakiri2022ProducedWaterReview,Adepitan2026HybridPW,Rajbongshi2024ProducedWaterReview,Ahmadun2009PWTechnologies}.

O processo de destilação por membrana é um processo de separação onde a força motriz é a diferença de pressão de vapor parcial decorrente da diferença de temperatura entre os dois lados da membrana. Tal membrana é hidrofóbica e microporosa, de modo a permitir apenas a passagem de vapor por meio dos poros, uma vez que a diferença de vapor parcial promove a evaporação do fluido no lado quente do módulo~\cite{Olatunji2018MDModeling,Yan2021MDWastewater,Kaur2023MDOverview,Adewole2022MDDesalination}.

O mecanismo para a recuperação do vapor varia de acordo com a geometria do módulo empregado, podendo ocorrer internamente ou externamente ao módulo. Em módulos de destilação por membrana com espaçamento de ar (AGMD) a condensação do vapor ocorre em contato com uma chapa resfriada localizada no interior do módulo, de modo que o permeado não entra em contato com o líquido de resfriamento. Nestes módulos, o espaçamento de ar é o fator mais impactante na transferência de calor e massa, uma vez que auxilia a reduzir a perda de calor por condução~\cite{Olatunji2018MDModeling,Yan2021MDWastewater,Kaur2023MDOverview,Adewole2022MDDesalination}. 

%A confiabilidade de um equipamento aplicado na industria está atrelada a capacidade de operar de acordo com os parâmetros de processo para o qual foi projetado sem apresentar falhas; A mantenabilidade de um equipamento reflete a capacidade e a facilidade da realização de reparos para restaurar a capacidade original de operação de um equipamento. Com base nesses conceitos, o conhecimento, o monitoramento e a compreensão dos modos de falha associados aos módulos de destilação por membranas tornam-se fundamentais para o adequado dimensionamento de seus componentes~\cite{smith2011}.

%Entre os possíveis modos de falha em módulos de destilação por membrana, destacam-se o molhamento (\textit{wetting}) e a incrustação (\textit{fouling}), fenômenos que exercem impacto direto na qualidade do permeado e na redução do fluxo ao longo do tempo. 

%O molhamento é um fenômeno onde o líquido a ser tratado permeia os poros da membrana hidrofóbica. Este fenômeno está ligado as propriedades morfológicas da membrana, as características do fluido e os parâmetros do processo. No que diz respeito a membrana, é possível citar o diâmetro de poro, a rugosidade e o ângulo de contato; Para o líquido, a tensão superficial e a presença de surfactantes; Para os parâmetros de processo, a velocidade do escoamento, a diferença de pressão na superfície da membrana e a temperatura~\cite{Yao2020WettabilityMD}.

%A incrustação no processo de destilação por membrana decorre de adsorção, acumulo ou precipitação na superfície da membrana. Assim como molhamento, este fenômeno é influenciado pelas propriedades da membrana, características do fluido e parâmetros do processo. A incrustação trás como principal consequência a redução da permeabilidade e consequentemente a redução do fluxo de permeado~\cite{Alkhatib2021FoulingMD}.

%Deste modo, para o correto dimensionamento e projeto de sistemas de tratamento de água de produção, é fundamental compreender o impacto dos modos de falha citados anteriormente em membranas expostas aos contaminantes presentes na água de produção por um longo período de tempo, em vista de selecionar o material adequado para compor o equipamento.

%Ante ao exposto, a presente dissertação tem como objetivo analisar os efeitos temporais em membranas comerciais durante o tratamento de água de produção via destilação por membranas, visando identificar a morfologia da membrana mais resistente às condições operacionais e avaliar a influência do tempo de operação no desempenho do processo.

Diante do exposto, a presente dissertação tem como objetivo investigar os efeitos temporais sobre membranas comerciais empregadas no tratamento de água de produção via destilação por membranas, buscando identificar as características morfológicas associadas à maior resistência às condições operacionais e avaliar a influência do tempo de operação sobre o desempenho e a estabilidade do processo. 

Para isso, realizou-se a caracterização das membranas a serem utilizadas, através de porosimetria de fluxo capilar, microscopia digital e medições de ângulo de contato. Também realizou-se a caracterização da água de produção através de métodos físico-químicos, cromatográficos e espectroscópicos para determinação dos contaminantes presentes. Por fim, serão conduzidos experiemntos de longo prazo utilizando um módulo AGMD, por um período de até trinta dias ou até a ocorrência de falhas operacionais, como molhamento ou incrustação, de modo a avaliar a estabilidade operacional e os mecanismos de degradação do processo ao longo do tempo.

%Esta dissertação é composta por este capítulo introdutório, seguida por outros cinco capítulos. O segundo capítulo abrange a revisão bibliográfica necessária para a compreensão do tema, abordando os aspectos morfológicos das membranas, as técnicas de tratamento de água de produção e os mecanismos de transferência de calor e massa presentes no processo de destilação por membranas com espaçamento de ar (AGMD). O terceiro capítulo apresenta a metodologia empregada na realização dos experimentos, a estrutura da bancada experimental e os parâmetros utilizados. O quarto capítulo apresenta e discute os resultados obtidos com os experimentos realizados. O quinto capítulo expõe as conclusões do trabalho e as considerações finais decorrentes do estudo realizado. Por fim, o sexto capítulo apresenta sugestões de trabalhos futuros.

